\chapter{Implementation}\label{chap:implementation}

This chapter presents the prototype implementation of the lattice-based RISC-V
ZkVM architecture described in the Chapter~\ref{chap:design}.
The implementation covers:

\begin{itemize}
	\item A subset of rv32imac instructions sufficient to run Nth Fibonacci and validate an Ethereum block with REVM.
	\item CCS constraints for these instructions.
	\item CCS constraints for the Poseidon2 IVC step commitment.
	\item CCS constraints for NIFS verifier.
\end{itemize}

Regarding the choice of folding scheme, this work uses the NetherMind's implementation of LatticeFold
protocol rather than its successor LatticeFold+, as the latter's library
implementation was not available at the time of system design and development.

\input{content/implementation/tech-stack.tex}

\input{content/implementation/crypto.tex}

\input{content/implementation/riscv-emul.tex}

\section{Circuit synthesis}

The constraints defining the validity of the VM execution and the IVC chain
are synthesized into a CCS.

\subsection{Poseidon2 constraint example}

This section illustrates the construction of CCS constraints for Poseidon2
hash \cite{Poseidon2}. Poseidon2 permutation in Goldilocks field on 12 elements (rate is 12,
with capacity 4, total width then being 16) consists of:

\begin{enumerate}
	\item Applying a maximum distance separable MDS matrix.
	\item 4 external initial rounds a 7th degree S-box on all elements.
	\item 22 internal rounds with a partial S-box applied only to the first element.
	\item 4 external terminal rounds a 7th degree S-box on all elements.
\end{enumerate}

\paragraph{Initial MDS application.}

The first operation applies the MDS matrix to the absorbed state. The input 12 elements
and 4 zeroes must be in the witness. Here this input will be denoted as $s = (s_0, s_1, \ldots, s_{11}, 0, 0, 0, 0)$.
Then the output of the MDS application is captured into the witness, denoted as $s_{after\_mds}$.
The constraint enforces

\begin{equation}
	s_{after\_mds\_i} - s'_i = 0
\end{equation}

where $i \in (0, 1, \ldots, 15)$ and $s' = \texttt{MDS} \cdot s$.

\paragraph{External initial rounds.}

Next are external rounds where a round constant is added and S-box
is applied to all state elements. There are 4 external initial rounds, after each
the state of 16 elements must be captured (denoted as
$s_{\text{ext\_init\_r\_i}}; \quad r \in (0, 1, 2, 3); \quad i \in (0, 1, \ldots, 15)$).
Then in circuit, for each round there is a constraint:

\begin{equation}
	\texttt{MDS}^{-1} \cdot s_{\text{ext\_init\_r\_i}} - (s_{\text{in\_r\_i}} + c_{\text{ext\_init\_r\_i}})^7 = 0
\end{equation}

where $c_{\text{ext\_init\_r\_i}}$ is the specific constant at index $i$ for round
$r$, and the $s_{\text{in\_r\_i}}$ is the input to the round, in first one, it
is the $s_{after\_mds}$, otherwise it is the output of previous round.

\paragraph{Internal rounds.}

Following the external initial rounds, the permutation executes 22 partial
internal rounds where the S-box is applied only to the first state element.
The state after each round is captured into witness, denoted as $s_{\text{inter\_r\_i}}; \quad r \in (0, 1, \ldots, 21); \quad i \in (0, 1, \ldots, 15)$.
Each internal round computes for each first element:

\begin{equation}
	M_I^{-1} \cdot s_{\text{inter\_r\_0}} - (s_{\text{in\_r\_0}} + c_{\text{inter\_0}})^7 = 0
\end{equation}

For other than first element ($i \in (1, 2, \ldots, 15)$ it does:

\begin{equation}
	M_I^{-1} \cdot s_{\text{inter\_r\_i}} - s_{\text{in\_r\_i}} = 0
\end{equation}

The $s_{\text{in\_0}}$ is the output of the last external initial round, other
$s_{\text{in\_i}}$ are outputs of the previous internal rounds.

\paragraph{External terminal rounds.}

The permutation concludes with 4 external terminal rounds, mirroring the
external initial rounds but using terminal round constants and different inputs
$s_{\text{ext\_term\_r\_i}}; \quad r \in (0, 1, 2, 3); \quad i \in (0, 1, \ldots, 15)$):

\begin{equation}
	\texttt{MDS}^{-1} \cdot s_{\text{ext\_term\_r\_i}} - (s_{\text{in\_r\_i}} + c_{\text{ext\_term\_r\_i}})^7 = 0
\end{equation}


where $c_{\text{ext\_term\_r\_i}}$ is the specific constant at index $i$ for round
$r$, and the $s_{\text{in\_r\_i}}$ is the input to the round, in first one, it
is the last internal round, otherwise it is the output of previous external terminal round.

\paragraph{Output extraction.}

After the final external terminal round, the output digest is extracted
from first 4 elements of the state

\begin{equation}
	h = (s_0, s_1, s_2, s_3)
\end{equation}

And compared to the claimed result.

\paragraph{Complexity analysis.}

For a single Poseidon2 sponge pass, the constraint system captures 388 witness
variables across the permutation computation. The system requires 44
matrices to encode the transformations, round constants, and S-box
operations. The number of matrices remains constant regardless of
the number of sponge passes, though their dimensions scale with the permutation
state size. The highest degree constraint is a 7th degree multiset constraint
corresponding to the Goldilocks S-box, and this degree does not increase with
additional sponge passes.

\subsection{Self-aware NIFS constraints}

Some constraints needed for the in-circuit NIFS verifier are not local in the
sense that they do not only depend on a fixed small set of witness variables,
but on the final structure of the whole CCS. The most notable example is the
linearization inner-product check, which verifies the relation

\begin{equation}
	\texttt{inner} = \sum_i c_i \cdot \prod_{j \in S_i} u_j,
\end{equation}

where the multisets $S_i$ are precisely the multisets of the CCS currently being
built. Therefore, this part of the verifier is self-aware. It must inspect the
complete list of already created multisets and generate constraints from them.

In the implementation this means that such constraints cannot be synthesized in
a single pass. The builder first creates all ordinary VM, Poseidon2,
decomposition, and folding constraints. Only after the full multiset structure
is known are the matrices for the linearization inner-product check
preallocated. Afterwards, in a second pass, these preallocated matrices are
filled row-by-row with the concrete indices corresponding to the final CCS
multisets. This is the reason why the folding related part of the builder must
ordered carefully, and why the constants describing the final CCS dimensions
(such as the number of matrices, multisets, and $\log_2(m)$) must be kept in
sync with the generated system.


\section{Implementation optimizations}

Besides the direct translation of algebraic relations into CCS constraints,
two optimizations were done in order to keep the ZkVM runnable on local
consumer-grade hardware as long as possible in development.

\subsection{\texttt{arkworks-rs} upgrade}

The LatticeFold library, and its dependencies, use the \texttt{arkworks-rs}\footnote{\url{https://github.com/arkworks-rs}}
version \texttt{v0.4.x} libraries for low level cryptography primitives. New versions
of these foundational libraries were released with optimizations and improvements.
In order to update and gain these benefits, the local fork of \texttt{latticefold}
was migrated to new version along with forking its dependencies
\texttt{cyclotomic-rings}, \texttt{stark-rings} and also migrating them to new versions.

\subsection{Reusing matrix structures}

Another optimization is that the builder groups together constraints
that share the same algebraic shape, even if they correspond to different parts
of the verifier. Instead of creating completely separate matrices and multisets
for every single row, the implementation often allocates one small family of
matrices and reuses them across many rows. For example, all product constraints
of the form "left input times right input equals intermediate product" can be
stored in the same pair of matrices and one linear
matrix for the output product. Likewise, entire families of recomposition
constraints are accumulated into shared linear matrices, differing only in the
rows that are populated. This reduces CCS size and avoids unnecessary growth in
the number of matrices and multisets.


\section{High level overview}

Figure~\ref{fig:execution_flow_overview} illustrates the complete execution flow. The
setup phase prepares the VM by loading the ELF binary and constructing the CCS
constraint system. Within the execution loop, four main modules work together:
the \texttt{riscvm} module handles instruction execution and trace generation,
the \texttt{commitments} module creates Poseidon2 hashes for state components,
the \texttt{ivc} module arithmetizes the trace into the CCS witness vector,
and the \texttt{main} module performs the folding operations that
compress each step into the accumulator. The dashed line on the right shows the
recursive nature of the computation, where the \texttt{IVCStepOutput} from step
$i$ becomes the \texttt{IVCStepInput} for step $i+1$.

\begin{tikzFigure}
	\begin{tikzpicture}[
			node distance=1.8cm and 1.0cm,
			font=\sffamily\small,
			% --- Styles ---
			% Functions / Logic
			func/.style={
					rectangle,
					draw=red!80!black,
					fill=red!10,
					thick,
					minimum height=0.7cm,
					minimum width=4.5cm,
					align=center,
					rounded corners=2pt
				},
			% Data Structures
			data/.style={
					rectangle,
					draw=blue!80!black,
					fill=blue!5,
					thick,
					minimum height=0.6cm,
					minimum width=4cm,
					align=center,
					dashed,
					font=\sffamily\footnotesize
				},
			% Module Containers
			module/.style={
					draw=gray!40,
					dashed,
					inner sep=0.5cm,
					rounded corners=5pt,
					fill=white,
					fill opacity=0.0
				},
			% Module Label
			modlabel/.style={
					font=\bfseries\scriptsize,
					text=gray!60!black,
					anchor=north east
				},
			% Lines
			line/.style={
					draw,
					-Latex,
					thick,
					gray!70!black
				}
		]

		% --- 1. SETUP PHASE (Top) ---
		\node[func] (setup_vm) {\texttt{VM::load\_elf}};
		\node[func, right=0.5cm of setup_vm] (setup_ccs) {\texttt{CCSBuilder::create\_riscv\_ccs}};

		% --- 2. THE LOOP  ---

		% --- Module: riscvm ---
		\node[func, below=1.8cm of setup_vm, xshift=2.5cm] (vm_exec) {\texttt{VM::fetch\_execute}};
		\node[data, below=0.3cm of vm_exec] (trace) {\texttt{ExecutionTrace}};

		% --- Module: commitments ---
		\node[func, below=1.5cm of trace] (committer) {
			\textbf{ZkVmCommitter methods:} \\
			ivc\_step\_comm,\\
			state\_i\_comm,\\
			acc\_comm,\\
			vm\_regs\_comm,\\
			vm\_mem\_comm
		};
		\node[data, below=0.3cm of committer] (comms) {Commitments represented by\\\texttt{GoldilocksComm}\\(4 Goldilocks field elements)};

		% --- Module: ivc ---
		\node[data, below=1.5cm of comms] (ivc_input) {\textbf{\texttt{IVCStepInput}} \\ (Prev Acc + Trace + Comms)};
		\node[func, below=0.3cm of ivc_input] (arith) {\texttt{arithmetize(input, layout)}};
		\node[data, below=0.3cm of arith] (z_vec) {\texttt{Vec<GoldilocksRingNTT>}\\($z$ vector)};

		% --- Module: main ---
		\node[func, below=1.5cm of z_vec] (commit_wit) {\texttt{commit(z, ccs, scheme)}};
		\node[data, below=0.3cm of commit_wit] (cccs) {\texttt{CCCS} ($cm_i$) + \texttt{Witness} ($w_i$)};

		\node[func, below=0.3cm of cccs] (nifs) {\texttt{fold(ccs, scheme, acc, w\_acc, cm\_i, w\_i)}};
		\node[data, below=0.3cm of nifs] (output) {\textbf{\texttt{IVCStepOutput}} \\ (New Accumulator + Proof)};

		% --- 3. MODULE BOUNDARIES ---
		\begin{scope}[on background layer]
			% VM Box
			\node[module, fit=(vm_exec) (trace)] (box_riscv) {};
			\node[modlabel] at (box_riscv.north east) {\texttt{mod riscvm}};

			% Comm Box
			\node[module, fit=(committer) (comms)] (box_comm) {};
			\node[modlabel] at (box_comm.north east) {\texttt{mod commitments}};

			% IVC Box
			\node[module, fit=(ivc_input) (arith) (z_vec)] (box_ivc) {};
			\node[modlabel] at (box_ivc.north east) {\texttt{mod ivc}};

			% Latticefold Box
			\node[module, fit=(commit_wit) (cccs) (nifs) (output)] (box_fold) {};
			\node[modlabel] at (box_fold.north east) {\texttt{crate latticefold}};

			% Main Loop Container
			\node[draw=gray!80, thick, inner sep=25pt, fit=(box_riscv) (box_fold), label={[anchor=north west, font=\bfseries]north west:Execution Loop (\texttt{vm.run})}] (main_loop) {};
		\end{scope}

		% --- 4. CONNECTIONS ---

		% Init -> Loop
		% Manual orthogonal lines
		\draw[line] (setup_vm.south east) -- (vm_exec.north);
		\draw[line] (setup_ccs.south east) -- ++(0,-4.5) -- (commit_wit.east);

		% Flow down
		\draw[line] (vm_exec) -- (trace);
		\draw[line] (trace) -- (committer);
		\draw[line] (committer) -- (comms);
		\draw[line] (comms) -- (ivc_input);
		\draw[line] (ivc_input) -- (arith);
		\draw[line] (arith) -- (z_vec);
		\draw[line] (z_vec) -- (commit_wit);
		\draw[line] (commit_wit) -- (cccs);
		\draw[line] (cccs) -- (nifs);
		\draw[line] (nifs) -- (output);

		\draw[line, dashed] (output.east) -- ++(2.9, 0)
		-- ++(0, 7.3)
		-- (ivc_input.east);

		% Label for loop
		\node[font=\scriptsize\bfseries, text=gray!80!black, rotate=90] at ($(output.east) + (2.7, 3.5)$) {Accumulator ($i \to i+1$)};
	\end{tikzpicture}
	\caption{High-level overview of the implemented execution flow.}
	\label{fig:execution_flow_overview}
\end{tikzFigure}

\pagebreak

\subsection{Toolchain summary}

The implementation is built with the following software stack:

\begin{itemize}
	\item \textbf{Rust:} \texttt{nightly-2025-08-19}, edition 2024.
	\item \textbf{Targets:} native host execution for the prover and \texttt{riscv32imac-unknown-none-elf} for guest programs.
	\item \textbf{LatticeFold:} a fork of Nethermind's \texttt{latticefold} implementation, used for Ajtai commitments, decomposition, linearization, and NIFS folding.
	\item \textbf{Plonky3 core crates:} \texttt{p3-field}, \texttt{p3-goldilocks}, \texttt{p3-poseidon2}, \texttt{p3-merkle-tree}, \texttt{p3-commit}, \texttt{p3-matrix}, \texttt{p3-mds}, \texttt{p3-challenger}, and \texttt{p3-symmetric}.
	\item \textbf{REVM:} version \texttt{36.0.0}, used inside the \texttt{evm} guest for replaying Ethereum block execution.
\end{itemize}

\subsection{Instruction coverage}

Emulator doesn't have a full \texttt{rv32imac} instruction coverage, just
enough for running a Fibonacci program and proving a block with REVM.
And also the zkVM doesn't constraint all the supported instructions.

\paragraph{Implemented in the VM emulator:}

\begin{itemize}
	\item \textbf{RV32I control flow:} \texttt{LUI}, \texttt{AUIPC}, \texttt{JAL}, \texttt{JALR}, \texttt{BEQ}, \texttt{BNE}, \texttt{BLT}, \texttt{BGE}, \texttt{BLTU}, \texttt{BGEU}.
	\item \textbf{RV32I memory:} \texttt{LB}, \texttt{LH}, \texttt{LW}, \texttt{LBU}, \texttt{LHU}, \texttt{SB}, \texttt{SH}, \texttt{SW}.
	\item \textbf{RV32I integer arithmetic and logic:} \texttt{ADDI}, \texttt{SLTI}, \texttt{SLTIU}, \texttt{XORI}, \texttt{ORI}, \texttt{ANDI}, \texttt{SLLI}, \texttt{SRLI}, \texttt{SRAI}, \texttt{ADD}, \texttt{SUB}, \texttt{SLL}, \texttt{SLT}, \texttt{SLTU}, \texttt{XOR}, \texttt{SRL}, \texttt{SRA}, \texttt{OR}, \texttt{AND}.
	\item \textbf{RV32M:} \texttt{MUL}, \texttt{MULHU}, \texttt{DIVU}, \texttt{REMU}.
	\item \textbf{RV32A:} \texttt{LR.W}, \texttt{SC.W}, \texttt{AMOADD.W}.
	\item \textbf{Other:} \texttt{FENCE}, \texttt{ECALL}. The local compressed decoder also contains an explicit fallback for \texttt{C.SUB}, \texttt{C.XOR}, \texttt{C.OR}, and \texttt{C.AND}, which are translated to their canonical RV32I forms.
\end{itemize}

\paragraph{Unsupported:}

\begin{itemize}
	\item \textbf{RV32I:} \texttt{EBREAK}.
	\item \textbf{RV32M:} \texttt{MULH}, \texttt{MULHSU}, \texttt{DIV}, \texttt{REM}.
	\item \textbf{RV32A:} \texttt{AMOSWAP.W}, \texttt{AMOXOR.W}, \texttt{AMOAND.W}, \texttt{AMOOR.W}.
	\item \textbf{RV32A:} \texttt{AMOMIN.W}, \texttt{AMOMAX.W}, \texttt{AMOMINU.W}, \texttt{AMOMAXU.W}.
\end{itemize}

\paragraph{Constrained inside the zkVM:}

\begin{itemize}
	\item \texttt{LUI}, \texttt{AUIPC}, \texttt{JAL}, \texttt{JALR}, \texttt{BNE}, \texttt{SW}, \texttt{ADDI}, and \texttt{ADD}.
\end{itemize}

\subsection{Parameter set}

The implementation fixes the following algebraic and cryptographic parameters:

\begin{itemize}
	\item \textbf{Base field:} the Goldilocks prime field with modulus
	      $p = 2^{64} - 2^{32} + 1 = 18446744069414584321$.
	\item \textbf{Cyclotomic ring:} coefficient representation over
	      $R_q = \mathbb{F}_q[X]/(X^{24} - X^{12} + 1)$, so the ring degree is 24.
	\item \textbf{CRT/NTT form:} the ring splits into 8 NTT components over
	      $\mathbb{F}_{q^3}$, i.e. extension degree 3 and \texttt{TAU} $= 24 / 8 = 3$.
	\item \textbf{Challenge set:} short lattice challenges use coefficients in the range $[-32, 32)$.
	\item \textbf{Ajtai/LatticeFold decomposition parameters:}
	      $B = 2^{15}$, $L = 5$, $B_{\text{small}} = 2$, and $K = 15$.
	\item \textbf{Commitment module rank:} $\kappa = 32$ ring elements.
	\item \textbf{Poseidon2 commitments:} output width 4, narrow state width/rate $(8, 4)$, wide state width/rate $(16, 12)$, full rounds 8, partial rounds 22, and S-box degree 7.
	\item \textbf{Hash security estimate:} the wide Poseidon2 sponge uses capacity 4 over a 64-bit field, which gives roughly 128-bit security.
	\item \textbf{Commitment matrix width:} the Ajtai commitment uses
	      $N = \texttt{w\_size} \cdot L = 19763 \cdot 5 = 98815$ columns.
\end{itemize}

\subsection{CCS dimensions}

The resulting instance has:

\begin{itemize}
	\item witness layout size $\texttt{w\_size} = 19763$,
	\item total $z$-vector size $n = 19768$,
	\item public input size $\ell = 4$,
	\item padded row count $m = 131072 = 2^{17}$,
	\item number of matrices $t = 125$,
	\item number of multisets $q = 52$,
	\item maximal multiset degree $d = 7$,
	\item row logarithm $s = 17$,
	\item column logarithm $s' = 15$.
\end{itemize}
